
Weight space embeddings enable model selection, editing, and analysis, yet most approaches flatten all parameters into a single vector, discarding architectural structure. This work investigates whether preserving network structure improves embedding quality for property prediction. Through systematic ablation on the CIFAR-10 Model Zoo (N = 61,335 models), layer-wise encoding---computing per-layer statistics (mean, standard deviation, minimum, maximum)---is compared against naïve flattening for accuracy prediction. Layer-wise encoding achieves Pearson correlation r = 0.547 versus r = 0.421 for Flatten+MLP, an improvement of $\Delta r$ = 0.126 (p < 0.001, t = 12.85). The effect holds across all five random seeds with 100\% consistency. Additional ablation steps involving Git Re-Basin alignment preprocessing could not be completed due to computational constraints at the 61K model scale. The CIFAR-10 Model Zoo is validated as a viable benchmark ($\sigma$ = 15.62\% accuracy variance). Results suggest that preserving layer boundaries provides substantial benefit for weight embedding quality.
